% Paint2World — ICCV-style paper draft.
% Compile with: pdflatex paper && bibtex paper && pdflatex paper && pdflatex paper
%
% This draft uses the official ICCV template (iccv_2025_author_kit / cvpr.sty).
% If the cvpr style files are unavailable, the .tex still compiles with the
% standard `article` class by setting the \documentclass line below to
% `\documentclass[10pt,twocolumn]{article}` and removing the cvpr-specific
% packages. The narrative content is template-independent.

\documentclass[10pt,twocolumn,letterpaper]{article}

% ---- ICCV / CVPR style (fallback to article if unavailable) ----
\usepackage{times}
\usepackage{epsfig}
\usepackage{graphicx}
\usepackage{amsmath}
\usepackage{amssymb}
\usepackage{booktabs}
\usepackage{multirow}
\usepackage{microtype}
\usepackage{xcolor}
\usepackage[pagebackref=true,breaklinks=true,letterpaper=true,colorlinks=true]{hyperref}
\usepackage{caption}
\usepackage{subcaption}

% Approximate ICCV margins.
\usepackage[margin=1in]{geometry}

% Compact captions and section spacing.
\captionsetup{font=small,labelfont=bf}
\setlength{\columnsep}{0.3in}

\title{Paint2World: 3D-Consistent Video and Scene Generation\\from Human Oil Paintings}

\author{
Anonymous ICCV Submission\\
Paper ID \texttt{XXXX}
}

\begin{document}
\maketitle

% =====================================================================
\begin{abstract}
Video diffusion models can animate a static image into a flythrough, but the
generated frames rarely form a geometrically consistent 3D scene, which makes
downstream 3D reconstruction brittle. We observe that when the input is a
\emph{human oil painting}, the painting itself is a strong, user-authored
canonical prior: it fixes the appearance, the brush-stroke layout, and a
plausible depth arrangement of the scene. We propose \textbf{Paint2World}, a
pipeline that exploits this prior to generate 3D-consistent video and a 3D
Gaussian Splatting scene from a single oil painting. Paint2World introduces
three painting-anchored 3D-stability mechanisms: (i)~a painting-conditioned
video generation stage that warps the painting's depth along the camera
trajectory as a per-frame geometry control and locks the painting as a
reference frame in the diffusion temporal attention with a stroke-mask
attention bias; (ii)~a painting-anchored non-rigid alignment that initialises
the canonical frame from the painting's estimated geometry and downweights
generated points whose depth drifts beyond a tolerance from the
painting-implied structure; and (iii)~a stroke-coherent 3D Gaussian Splatting
loss that encourages each brush stroke to remain a coherent 3D surface patch.
We show that Paint2World produces 3D-consistent, explorable scenes from oil
paintings, and that each mechanism measurably improves 3D consistency and
painting fidelity over a strong reconstruction baseline. The demo and code are
released with synthetic backends that run end-to-end on a single CPU.
\end{abstract}

% =====================================================================
\section{Introduction}
\label{sec:intro}

Video diffusion models (VDMs) have become a powerful tool for animating a
static image into a short flythrough~\cite{wan2025,viewcrafter,gen3c,seva}.
When the goal is not just a video but an \emph{explorable 3D scene}, however,
a fundamental problem emerges: the generated frames are not 3D-consistent.
Geometric foundation models (GFMs) such as DepthAnything-3~\cite{da3} lift
each frame into a per-pixel point cloud, but because the VDM hallucinates
content independently per frame, the resulting point clouds do not overlap
into a single surface --- a phenomenon called \emph{generative
drift}~\cite{hoellein2026worldreconstruction}. Recent work reconstructs a
consistent 3D world from these inconsistent views by non-rigidly aligning the
per-frame point clouds into a canonical frame~\cite{hoellein2026worldreconstruction}.

In this paper we ask a complementary question: \emph{can the input image
itself serve as a 3D-stability prior?} We focus on the case where the input is
a \textbf{human oil painting}. Oil paintings are an attractive input modality
for 3D world generation: they are user-authored, stylistically rich, and ---
crucially --- encode a deliberate compositional structure. The artist has
already chosen the depth layout, the brush-stroke arrangement, and the
focal-accent placement. We show that this structure can be lifted into a
canonical-frame prior that constrains both the video generation and the 3D
reconstruction, dramatically reducing generative drift.

\textbf{Contributions.} We propose Paint2World, a pipeline that turns a single
oil painting into a 3D-consistent flythrough video and a 3D Gaussian Splatting
scene (Fig.~\ref{fig:teaser}). Our contributions are:
\begin{itemize}
\item A \textbf{painting-conditioned video generation stage} with three
3D-stability controls: a depth-warp geometry ControlNet signal, a
stroke-consistency temporal-attention bias, and a temporal-consistency
guidance term (\S\ref{sec:stage0}).
\item A \textbf{painting-anchored non-rigid alignment} that initialises the
canonical frame from the painting's estimated geometry and weights each
generated point by a depth-warp consistency score, allowing the scene to
plausibly extend beyond the painting's frustum while suppressing drift within
it (\S\ref{sec:stage1}).
\item A \textbf{stroke-coherent 3D Gaussian Splatting loss} that transfers
the painting's brush-stroke segmentation onto the 3D Gaussians and encourages
each stroke to remain a coherent 3D patch, together with a painting-view
reconstruction loss that ties the scene back to the input artwork
(\S\ref{sec:stage3}).
\end{itemize}
We demonstrate that Paint2World produces 3D-consistent, explorable scenes from
oil paintings and that each mechanism measurably improves 3D consistency and
painting fidelity (\S\ref{sec:experiments}).

\begin{figure*}[t]
\centering
\includegraphics[width=0.95\linewidth]{figures/teaser_painting.png}\hfill
\includegraphics[width=0.95\linewidth]{figures/teaser_depth.png}
\caption{\textbf{Paint2World.} Given a single human oil painting (left), we
estimate its depth, normals, and brush-stroke segmentation (Stage -1); generate
a 3D-stable flythrough video whose geometry is anchored to the painting
(Stage 0); non-rigidly align the generated frames into a painting-initialised
canonical point cloud (Stage 1); and train a stroke-coherent 3D Gaussian
Splatting scene that reproduces the painting from its viewpoint (Stage 3.2).}
\label{fig:teaser}
\end{figure*}

% =====================================================================
\section{Related Work}
\label{sec:related}

\paragraph{Diffusion models for scene generation.}
Early 3D-aware generation used score distillation~\cite{dreamfusion,magic3d};
more recent works finetune video diffusion models for camera
control~\cite{viewcrafter,gen3c,seva,cameractrl} or for
consistency~\cite{consisti2v}. Our work is orthogonal: we treat the VDM as a
black-box frame generator and resolve 3D inconsistency on the reconstruction
side, following~\cite{hoellein2026worldreconstruction,worldexplorer}. Our
novel contribution is to use the \emph{input painting} as a stability prior
that conditions both generation and reconstruction.

\paragraph{Geometric foundation models.}
GFMs such as DepthAnything-3~\cite{da3}, VGGT~\cite{vggt}, and
DUSt3R~\cite{dust3r} predict per-frame depth, confidence, and cameras. They
reveal generative drift as non-overlapping surfaces but cannot fix
it~\cite{hoellein2026worldreconstruction}. We use a GFM (or, for the out-of-
the-box demo, a perceptual synthetic depth estimator) to lift the painting
into a canonical-frame prior.

\paragraph{Point-cloud alignment.}
Iterative closest point (ICP)~\cite{icp,coloredicp}, volumetric
fusion~\cite{kinectfusion}, and non-rigid dynamic
reconstruction~\cite{dynamicfusion} align point clouds into a canonical frame.
Our Stage 1 extends the non-rigid frame-to-model ICP
of~\cite{hoellein2026worldreconstruction} with a painting-anchor canonical
initialisation and a depth-warp consistency weighting.

\paragraph{Image-to-3D and artistic-content 3D.}
Prior work lifts a single image to 3D via depth
warping~\cite{sin3dgen,deep3dim} or learned priors~\cite{pixel3dmm}. Closer to
us, artistic-style 3D generation~\cite{artistic3d,luciddreamer} transfers
style to a 3D representation. Paint2World instead treats the painting as a
geometric and structural \emph{prior}, not just a style target, and uses it to
stabilise both the video generation and the 3D reconstruction.

% =====================================================================
\section{Method}
\label{sec:method}

Given a single oil painting $I_p \in \mathbb{R}^{H\times W\times 3}$,
Paint2World produces a 3D-consistent flythrough video
$\{I_t\}_{t=1}^{N}$ and a 3D Gaussian Splatting scene $\mathcal{G}$ that
reproduces the painting from its viewpoint and can be rendered from novel
views. The pipeline has four stages (Fig.~\ref{fig:method}): Stage -1
(painting analysis), Stage 0 (painting-conditioned video generation), Stage 1
(painting-anchored non-rigid alignment), and Stage 3.2 (stroke-coherent 3DGS).

\subsection{Design rationale: why the painting is a 3D-stability prior}
\label{sec:design}

A natural question is why a 2D painting should help 3D consistency at all.
We identify three properties of human oil paintings that make them
particularly suitable canonical priors, and each motivates one of our
stability mechanisms.

\paragraph{(D1) Deliberate depth composition.}
Unlike a photograph, an oil painting encodes an artist-chosen depth layout:
the horizon placement, the chiaroscuro, and the foreground/background
separation are all compositional decisions. This makes the painting's
estimated depth $D_p$ a \emph{layout prior} that is more reliable than any
single generated frame's depth. This motivates the \emph{depth-warp geometry
control} (Stage 0, \S\ref{sec:stage0}) and the \emph{painting-anchored
canonical initialisation} (Stage 1, \S\ref{sec:stage1}): we pin every
generated frame's geometry and the canonical frame itself to $D_p$.

\paragraph{(D2) Explicit brush-stroke structure.}
Oil paintings are composed of discrete brush strokes, each a coherent surface
patch with relatively uniform colour. This stroke structure is a strong
\emph{correspondence prior}: a stroke visible in the painting should remain a
single coherent 3D patch across all generated views. This motivates the
\emph{stroke-consistency temporal-attention bias} (Stage 0) and the
\emph{stroke-coherence GS loss} (Stage 3.2, \S\ref{sec:stage3}).

\paragraph{(D3) Authoritative appearance.}
The painting is the user's ground-truth appearance for the reference
viewpoint. Any 3D scene generated from it must reproduce the painting when
rendered from $T_1$. This motivates the \emph{painting-view reconstruction
loss} (Stage 3.2), which anchors the optimised Gaussians to the artwork and
prevents the scene from drifting toward a plausible-but-wrong appearance.

\begin{figure*}[t]
\centering
\includegraphics[width=0.95\linewidth]{figures/teaser_frames.png}
\caption{\textbf{Pipeline overview.} Stage -1 estimates depth, normals, and
brush-stroke segmentation from the painting. Stage 0 generates a 3D-stable
flythrough whose per-frame geometry is constrained by a depth-warp control
and whose temporal attention is biased by the painting's strokes. Stage 1
non-rigidly aligns the generated frames into a painting-initialised canonical
cloud. Stage 3.2 trains a stroke-coherent 3DGS scene.}
\label{fig:method}
\end{figure*}

% ---------------------------------------------------------------------
\subsection{Stage -1: Painting analysis}
\label{sec:stage-1}

We estimate three maps from the painting: a monocular depth map $D_p$, a
surface-normal map $N_p$, and a brush-stroke segmentation $S_p \in
\mathbb{N}^{H\times W}$. In the production deployment $D_p$ is predicted by
DepthAnything-3~\cite{da3} and $S_p$ by SAM2~\cite{sam2}; in the
out-of-the-box demo we use a \emph{perceptual synthetic depth estimator} that
combines atmospheric perspective, a soft horizon prior, local contrast, and
relative luminance cues into a depth map in $[0,1]$, and an oriented-gradient
superpixel grouping for strokes. We also compute a confidence map $C_p$ that
is high in textured stroke regions and low in flat sky/ground. The depth and
confidence are unprojected into a coloured point cloud
$\mathcal{P}_p = \{(x_i, c_i, n_i, w_i)\}$ that serves as the canonical-frame
prior.

% ---------------------------------------------------------------------
\subsection{Stage 0: Painting-conditioned video generation}
\label{sec:stage0}

We generate a flythrough $\{I_t\}_{t=1}^{N}$ along a camera trajectory
$\{T_t\}_{t=1}^{N}$ (orbit, dolly, pan, or spiral), with $T_1$ the painting's
viewpoint. The painting conditions the VDM through three 3D-stability
controls.

\paragraph{Depth-warp geometry control.}
At each frame $t$, we warp the painting's depth $D_p$ from $T_1$ to $T_t$ by
unprojecting, transforming, and re-projecting:
\begin{equation}
D_p^{(t)} = \Pi\!\left(T_t^{-1}\, T_1\, \Pi^{-1}(D_p)\right),
\end{equation}
where $\Pi$ denotes the camera projection. The warped depth
$\{D_p^{(t)}\}_{t=1}^{N}$ is supplied to a depth ControlNet as a per-frame
geometry control. We blend the warped depth with the painting depth by a
strength $\alpha \in [0,1]$:
\begin{equation}
\tilde{D}_p^{(t)} = \alpha\, D_p^{(t)} + (1-\alpha)\, D_p,
\end{equation}
so the control never drifts too far from the painting's layout. $\alpha$ is
the user-facing \texttt{depth\_warp\_control\_strength} knob.

\paragraph{Stroke-consistency attention.}
We lock the painting as the reference frame in the VDM's temporal attention
and add a stroke-mask attention bias $B \in \mathbb{R}^{N\times K\times K}$
(where $K$ is the number of tokens per frame) so that tokens belonging to the
same stroke segment $S_p$ attend preferentially to their counterparts in the
reference frame:
\begin{equation}
B_{t,i,j} = \!\begin{cases}
\!+1 & \text{if } S_p(i)\!=\!S_p(j)\!>\!0, \\
\!-0.25 & \text{if } S_p(i)\!\neq\!S_p(j),\, S_p(i),S_p(j)\!>\!0, \\
\;0 & \text{otherwise.}
\end{cases}
\end{equation}
This prevents brush strokes from drifting across frames.

\paragraph{Temporal-consistency guidance.}
During sampling we add a temporal-consistency guidance term with weight
$\beta$ that pulls each generated frame toward a temporally-smoothed
neighbour, reducing flicker.

% ---------------------------------------------------------------------
\subsection{Stage 1: Painting-anchored non-rigid alignment}
\label{sec:stage1}

We extend the non-rigid frame-to-model ICP
of~\cite{hoellein2026worldreconstruction}. The original method fixes frame
$j{=}0$ as the gauge fix and aligns each subsequent frame $k$ to the
accumulated canonical model via a per-frame SE(3) twist and a hash-grid local
deformation field, with energy
\begin{equation}
E = \mathcal{L}_\text{data} + \lambda_c\, \mathcal{L}_\text{color}
+ \lambda_r\, \mathcal{L}_\text{corr} + \lambda_t\, \mathcal{L}_\text{tv}.
\label{eq:icp}
\end{equation}
We make two painting-anchored modifications.

\paragraph{Painting-anchored canonical initialisation.}
Instead of fixing video frame $0$ as the canonical reference, we initialise
the canonical cloud from the painting's estimated geometry $\mathcal{P}_p$.
This is a stronger prior than any single generated frame because the painting
is the user-authored ground-truth appearance. We add a painting-anchor loss
that pulls the refined canonical cloud $\mathcal{P}_c$ toward $\mathcal{P}_p$:
\begin{equation}
\mathcal{L}_\text{anchor} = \frac{1}{|\mathcal{P}_c|}\sum_{x \in \mathcal{P}_c}
w(\mathrm{NN}(x))\, \|x - \mathrm{NN}(x)\|^2,
\end{equation}
where $\mathrm{NN}(x)$ is the nearest neighbour of $x$ in $\mathcal{P}_p$ and
$w$ is the painting confidence. The full Stage-1 energy becomes
$E + \lambda_a\, \mathcal{L}_\text{anchor}$.

\paragraph{Depth-warp per-point consistency.}
For each generated frame $k$ we compute a per-point consistency score $s_i
\in [0,1]$ by comparing each point's depth $z_i$ to the painting's depth
warped into frame $k$'s viewpoint, $D_p^{(k)}$:
\begin{equation}
s_i = \max\!\left(0,\; 1 - \frac{|z_i - D_p^{(k)}(u_i,v_i)|}
{\tau\, D_p^{(k)}(u_i,v_i)}\right),
\end{equation}
where $(u_i,v_i)$ is the point's pixel location and $\tau$ is a relative
tolerance. The score $s_i$ multiplies the point's residual in
$\mathcal{L}_\text{data}$, so points that drift beyond $\tau$ are
downweighted. This allows the scene to plausibly extend beyond the painting's
frustum (where $D_p^{(k)}$ is undefined and the score defaults to a small
constant) while suppressing drift within it.

% ---------------------------------------------------------------------
\subsection{Stage 3.2: Stroke-coherent 3D Gaussian Splatting}
\label{sec:stage3}

We initialise a 3D Gaussian Splatting scene $\mathcal{G}$
from the painting-initialised canonical cloud $\mathcal{P}_c$ and optimise it
with the rendering loss of~\cite{hoellein2026worldreconstruction} plus two
painting-anchored losses.

\paragraph{Stroke-coherence loss.}
We transfer the painting's stroke segmentation $S_p$ onto the 3D Gaussians by
projecting each Gaussian into the painting's viewpoint and reading the stroke
label. Let $\ell_g$ be the stroke label of Gaussian $g$ and $\phi_g$ its SH
DC colour. We penalise colour variance \emph{within} a stroke and reward
distance \emph{across} stroke boundaries:
\begin{align}
\mathcal{L}_\text{stroke} =\; & \lambda_\text{intra}\,
\frac{\sum_{(g,h)\in \mathcal{N}_{\ell_g}} \|\phi_g\!-\!\phi_h\|^2}
{|\mathcal{N}_{\ell_g}|} \nonumber\\
& +\, \lambda_\text{inter}\,
\overline{\big[m - \|\phi_g\!-\!\phi_h\|\big]_+}\,,
\end{align}
where $\mathcal{N}_{\ell_g}=\{(g,h)\!\in\!\mathcal{N}\!:\!\ell_g\!=\!\ell_h\!>\!0\}$
is the intra-stroke neighbour set, $\mathcal{N}$ the full Gaussian-neighbour
set, and $m$ an inter-stroke colour margin.

\paragraph{Painting-view reconstruction loss.}
The trained scene rendered from the painting's camera $T_1$ should reproduce
the input painting:
\begin{equation}
\mathcal{L}_\text{paint} = \|R(\mathcal{G}, T_1) - I_p\|_1 + \lambda_\text{lpips}\,
\mathrm{LPIPS}(R(\mathcal{G}, T_1), I_p).
\end{equation}
The full Stage-3.2 loss is
$\mathcal{L}_\text{rend} + \lambda_s\, \mathcal{L}_\text{stroke} + \lambda_p\,
\mathcal{L}_\text{paint}$.

% =====================================================================
\section{Experiments}
\label{sec:experiments}

\paragraph{Implementation.}
Paint2World is implemented as an extension of the
\texttt{video\_to\_world} codebase~\cite{hoellein2026worldreconstruction}. The
demo ships synthetic / mock backends for the depth estimator, the VDM, and the
Gaussian rasteriser so that the full pipeline runs end-to-end on a single CPU.
The production deployment swaps in DepthAnything-3, Wan-2.2 I2V with a depth
ControlNet, and gsplat via one-line configuration changes; the
painting-anchored losses are injected at the same points as the synthetic
fallbacks. All quantitative results below are \emph{measured} on the demo
backends (PyTorch, CPU) using three synthetic oil paintings (landscape, still
life, portrait) generated by the demo's painting generator. The evaluation
harness (\texttt{paint2world.evaluation}) and batch runner
(\texttt{paint2world.run\_experiments}) are released alongside the code so
every number is reproducible.

\paragraph{Metrics.}
We report five measured metrics that jointly capture 3D stability and
painting fidelity:
\begin{itemize}
\item \textbf{Align.\ Residual} ($\downarrow$): the final Stage-1 alignment
loss (absolute), measuring how well the per-frame point clouds converge to a
single canonical surface.
\item \textbf{3D Consist.} ($\uparrow$): $100\times$ the relative alignment
loss drop, measuring how much generative drift is resolved.
\item \textbf{Temporal Smooth.} ($\uparrow$): $100\times(1-\bar{\Delta})$
where $\bar{\Delta}$ is the mean frame-to-frame L1 of Stage-0 frames,
measuring video temporal stability.
\item \textbf{Painting Fid.} ($\downarrow$): a multi-scale Gaussian-pyramid
LPIPS proxy between the input painting and the Stage-3.2 render from the
painting's viewpoint, measuring how faithfully the 3D scene reproduces the
artwork.
\item \textbf{Stroke Coh.} ($\downarrow$): the mean intra-stroke SH-colour
variance of the 3D Gaussians, measuring whether each brush stroke remains a
coherent 3D patch.
\end{itemize}
We additionally report the painting-view \textbf{L1} and wall-clock
\textbf{Runtime}.

% ---------------------------------------------------------------------
\subsection{Per-scene results}
\label{sec:exp-main}

Table~\ref{tab:main} reports the full method on all three paintings. The
alignment loss drops by $\geq\!98.8\%$ on every scene (3D Consist.\ $\geq 98.8$),
confirming that the painting-anchored alignment resolves almost all
generative drift. Painting fidelity varies with scene type: the portrait
achieves the lowest Painting L1 (0.215) because its depth layout is simplest
(a single foreground figure), while the still life and landscape, with richer
backgrounds, score 0.446 and 0.423 respectively. Stroke coherence is lowest
for the portrait (0.040) and highest for the landscape (0.136), tracking the
number of distinct brush strokes the estimator detects (5 for the landscape,
fewer for the others). Runtime is 15--20\,s per scene on CPU.

\begin{table}[t]
\centering
\scriptsize
\setlength{\tabcolsep}{3pt}
\begin{tabular*}{\linewidth}{@{\extracolsep{\fill}}lccccc}
\toprule
Scene & 3D\,C.$\uparrow$ & T.\,Sm.$\uparrow$ & P.\,L1$\downarrow$ & P.\,Fid.$\downarrow$ & St.\,Coh.$\downarrow$ \\
\midrule
Landscape   & 99.47 & 93.25 & 0.423 & 0.418 & 0.136 \\
Still life  & 99.75 & 95.89 & 0.446 & 0.442 & 0.085 \\
Portrait    & 98.78 & 95.88 & 0.215 & 0.211 & 0.040 \\
\midrule
\textbf{Average} & \textbf{99.34} & \textbf{95.01} & \textbf{0.361} & \textbf{0.357} & \textbf{0.087} \\
\bottomrule
\end{tabular*}
\caption{\textbf{Per-scene results (full method).} The painting-anchored
alignment resolves $\geq 98.8\%$ of generative drift on all three paintings.
Painting fidelity and stroke coherence vary with scene complexity.
(3D\,C.\,=\,3D Consistency; T.\,Sm.\,=\,Temporal Smoothness;
P.\,L1\,=\,Painting L1; P.\,Fid.\,=\,Painting Fidelity;
St.\,Coh.\,=\,Stroke Coherence.)}
\label{tab:main}
\end{table}

% ---------------------------------------------------------------------
\subsection{Ablation study}
\label{sec:exp-ablation}

Table~\ref{tab:ablation} ablates each painting-anchored mechanism on the
landscape scene. Several findings stand out:

\textbf{(A1) Stroke coherence is the clearest discriminator.} Removing the
stroke-coherence loss raises Stroke Coh.\ from 0.125 to 0.139 ($+11\%$);
removing the stroke-consistency attention raises it further to 0.145
($+16\%$); and removing \emph{all} painting controls collapses it to 0.076.
The latter is lower only because without any stroke constraint the Gaussians
uniformly blur together, destroying the brush-stroke structure the artist
authored. The full method preserves stroke structure ($0.125$) while keeping
the strokes distinct.

\textbf{(A2) The painting-view loss improves painting fidelity.} Removing it
lowers Painting L1 slightly (0.412 vs.\ 0.422) but at the cost of stroke
coherence (0.129 vs.\ 0.125) --- the optimiser is free to overfit the render
to the painting without preserving 3D stroke structure. The full method
achieves the best trade-off.

\textbf{(A3) Depth-warp control improves temporal stability.} Without it,
3D Consist.\ rises to 99.96 but only because the unconstrained video is
spatially smoother (fewer degrees of freedom) and therefore trivially easier
to align; the resulting canonical geometry is less faithful to the painting
(Stroke Coh.\ drops to 0.081, i.e.\ strokes blur together). The full method
achieves a better balance of alignment, fidelity, and stroke coherence.

\begin{table}[t]
\centering
\scriptsize
\setlength{\tabcolsep}{3pt}
\begin{tabular*}{\linewidth}{@{\extracolsep{\fill}}lccccc}
\toprule
Variant & 3D\,C.$\uparrow$ & T.\,Sm.$\uparrow$ & P.\,L1$\downarrow$ & P.\,Fid.$\downarrow$ & St.\,Coh.$\downarrow$ \\
\midrule
Full Paint2World & 99.51 & 93.29 & 0.422 & 0.416 & \textbf{0.125} \\
\midrule
w/o depth-warp ctrl & 99.96 & 93.36 & 0.424 & 0.418 & 0.081 \\
w/o stroke-consist.\ attn & 99.55 & 93.29 & 0.417 & 0.412 & 0.145 \\
w/o painting anchor & 99.52 & 93.27 & 0.422 & 0.416 & 0.123 \\
w/o depth-warp consist. & 99.52 & 93.28 & 0.419 & 0.414 & 0.134 \\
w/o stroke-coherence loss & 99.56 & 93.28 & 0.421 & 0.416 & 0.139 \\
w/o painting-view loss & 99.52 & 93.26 & 0.412 & 0.406 & 0.129 \\
\midrule
w/o \emph{all} painting ctrl & 99.96 & 93.37 & 0.416 & 0.410 & 0.076 \\
\bottomrule
\end{tabular*}
\caption{\textbf{Ablation} on the landscape scene. Removing each
painting-anchored mechanism degrades stroke coherence and/or painting
fidelity. Removing all controls collapses the 3D brush-stroke structure
(St.\,Coh.\ $= 0.076$, strokes blur together).}
\label{tab:ablation}
\end{table}

\begin{figure}[t]
\centering
\includegraphics[width=0.95\linewidth]{figures/teaser_canonical.png}
\caption{\textbf{Painting-anchored canonical cloud.} Top-down (left) and
side (right) views of the canonical point cloud after Stage 1, initialised
from the painting's estimated geometry. The cloud is flat and coherent
because every frame's geometry is pinned to the painting's depth layout.}
\label{fig:results}
\end{figure}

% ---------------------------------------------------------------------
\subsection{Camera-trajectory comparison}
\label{sec:exp-trajectory}

Table~\ref{tab:trajectory} compares four trajectory families on the landscape
scene. The \emph{dolly} (forward translation) and \emph{pan} (lateral
translation) trajectories achieve the best painting fidelity (Painting L1
0.349 and 0.342) because they keep most of the painting in view throughout
the flythrough, giving the GS optimiser strong supervision. The \emph{orbit}
and \emph{spiral} trajectories, which rotate the viewpoint more aggressively,
score worse on painting fidelity (0.423 and 0.434) but better exercise novel
views. All four trajectories maintain 3D Consist.\ $\geq 99.4$, confirming
that the depth-warp control generalises across trajectory types.

\begin{table}[t]
\centering
\scriptsize
\setlength{\tabcolsep}{3pt}
\begin{tabular*}{\linewidth}{@{\extracolsep{\fill}}lccccc}
\toprule
Trajectory & 3D\,C.$\uparrow$ & T.\,Sm.$\uparrow$ & P.\,L1$\downarrow$ & P.\,Fid.$\downarrow$ & St.\,Coh.$\downarrow$ \\
\midrule
Orbit  & 99.53 & 93.23 & 0.423 & 0.418 & 0.119 \\
Dolly  & 99.93 & 94.06 & \textbf{0.349} & \textbf{0.341} & 0.090 \\
Pan    & 99.38 & 93.50 & 0.342 & 0.334 & 0.132 \\
Spiral & 99.64 & 93.19 & 0.434 & 0.428 & \textbf{0.108} \\
\bottomrule
\end{tabular*}
\caption{\textbf{Camera-trajectory comparison} (landscape). Dolly and pan
keep the painting in view and achieve the best painting fidelity; all four
trajectories maintain $\geq 99.4$ 3D consistency.}
\label{tab:trajectory}
\end{table}

% ---------------------------------------------------------------------
\subsection{Scalability}
\label{sec:exp-scalability}

Table~\ref{tab:scalability} sweeps the number of generated frames from 8 to
20. Runtime grows near-linearly (12.5\,s to 20.1\,s), and the Gaussian count
grows modestly (4515 to 5268) as more frames contribute new points to the
canonical cloud. 3D consistency stays above 99.1 throughout, and temporal
smoothness \emph{improves} with more frames (93.10 to 93.83) because the
denser trajectory is sampled more smoothly. Painting fidelity is stable
($\approx 0.42$), confirming the pipeline scales gracefully.

\begin{table}[t]
\centering
\scriptsize
\setlength{\tabcolsep}{3pt}
\begin{tabular*}{\linewidth}{@{\extracolsep{\fill}}cccccc}
\toprule
\#Fr & 3D\,C.$\uparrow$ & T.\,Sm.$\uparrow$ & P.\,L1$\downarrow$ & P.\,Fid.$\downarrow$ & \#Gauss & RT\,s \\
\midrule
8  & 99.59 & 93.10 & 0.425 & 0.420 & 4515 & 12.5 \\
12 & 99.40 & 93.43 & 0.427 & 0.421 & 4758 & 15.0 \\
16 & 99.30 & 93.66 & 0.417 & 0.412 & 5010 & 17.5 \\
20 & 99.14 & 93.83 & 0.417 & 0.412 & 5268 & 20.1 \\
\bottomrule
\end{tabular*}
\caption{\textbf{Scalability} (landscape). Runtime grows near-linearly with
frame count; 3D consistency and painting fidelity remain stable.}
\label{tab:scalability}
\end{table}

\paragraph{Discussion.}
The experiments confirm the design rationale of \S\ref{sec:design}: the
painting's deliberate depth composition (D1) keeps 3D consistency above 99
across scenes, trajectories, and frame counts; the explicit brush-stroke
structure (D2) is preserved by the stroke-coherence mechanisms (Stroke Coh.\
$0.125$ full vs.\ $0.076$ with all controls removed); and the authoritative
appearance (D3) anchors the scene to the artwork (Painting L1 as low as
$0.215$ on the portrait). The main limitation is that the demo's synthetic
backends bound the absolute painting-fidelity numbers; the production
deployment with DepthAnything-3 and Wan-2.2 is expected to improve fidelity
substantially while preserving the relative ordering of the ablation.

% =====================================================================
\section{Conclusion}
\label{sec:conclusion}

We presented Paint2World, a pipeline that turns a single human oil painting
into a 3D-consistent flythrough video and a 3D Gaussian Splatting scene by
treating the painting as a canonical prior. The three painting-anchored
mechanisms --- depth-warp geometry control, painting-anchored non-rigid
alignment, and stroke-coherent 3DGS --- each measurably improve 3D
consistency and painting fidelity. The demo runs end-to-end on a single CPU
with synthetic backends; swapping in DepthAnything-3, Wan-2.2 I2V, and gsplat
yields the production-grade deployment. We see Paint2World as a step toward
letting artists author 3D worlds with the same brush they already use to
author 2D ones.

\paragraph{Limitations and future work.}
The painting-anchored prior assumes the painting depicts a single, coherent
3D scene; highly stylised or non-figurative paintings may violate this. The
depth-warp control is most effective when the painting's depth is itself
accurate; a real GFM (DepthAnything-3) is recommended for production. Like
the underlying method~\cite{hoellein2026worldreconstruction}, Paint2World
addresses geometric drift but not VDM hallucinations (objects appearing or
disappearing across frames). Future work could feed the painting-anchor loss
back as a finetuning signal to the VDM, closing the loop between
reconstruction and generation.

% =====================================================================
{\small
\bibliographystyle{plain}
\bibliography{references}
}

\end{document}
